% =====================================================================
%  CV - Ibrahim Kakaev
%  Modèle « gouttière » : une colonne étroite à gauche pour les dates et
%  intitulés, le contenu sur tout le reste. Une seule colonne de lecture,
%  sans tableau ni zone de texte : compatible ATS.
%  Compilation : pdfLaTeX (compilateur par défaut d'Overleaf)
%  Photo       : importer "photo.PNG" à la racine du projet (optionnel).
% =====================================================================
\documentclass[10pt,a4paper]{article}

\usepackage[utf8]{inputenc}
\usepackage[T1]{fontenc}
\usepackage[sfdefault]{FiraSans}
\usepackage[a4paper,top=1.05cm,bottom=0.9cm,left=1.3cm,right=1.3cm]{geometry}
\usepackage{xcolor}
\usepackage{graphicx}
\usepackage{enumitem}
\usepackage{calc}
\usepackage[hidelinks]{hyperref}
\usepackage[letterspace=120]{microtype}

% Texte du PDF lisible par les ATS (accents, ligatures)
\ifdefined\pdfgentounicode
  \input{glyphtounicode}
  \pdfgentounicode=1
\fi

\pagestyle{empty}
\setlength{\parindent}{0pt}
\setlength{\parskip}{0pt}
\raggedright

% ---------- Noir et gris uniquement ----------
\definecolor{ink}{HTML}{161616}
\definecolor{soft}{HTML}{6A6A6A}
\definecolor{hair}{HTML}{CFCFCF}
\color{ink}

% ---------- Grille : largeur de la gouttière ----------
\newlength{\gut}\setlength{\gut}{2.55cm}

% Titre de section : barre épaisse, intitulé espacé, filet jusqu'au bord
\newcommand{\cvsection}[1]{%
  \vspace{7pt}%
  {\color{ink}\rule[-0.5pt]{2.4pt}{9pt}}\hspace{6pt}%
  {\small\bfseries\MakeUppercase{#1}}\hspace{8pt}%
  {\color{hair}\leaders\hrule height 3.4pt depth -2.9pt\hfill}\par
  \vspace{3.5pt}}

% Bloc à gouttière : l'étiquette (date, intitulé) à gauche, le contenu à droite
\newenvironment{grid}{%
  \begin{list}{}{\setlength{\leftmargin}{\gut}\setlength{\labelwidth}{\gut-8pt}%
    \setlength{\labelsep}{8pt}\setlength{\itemsep}{4pt}\setlength{\topsep}{0pt}%
    \setlength{\parsep}{0pt}\setlength{\partopsep}{0pt}%
    \renewcommand{\makelabel}[1]{\makebox[\labelwidth][l]{\small\color{soft}##1}}}}
  {\end{list}}

% Puces
\newenvironment{cvitems}{\begin{itemize}[leftmargin=0.95em,itemsep=0pt,topsep=1pt,parsep=0pt,partopsep=0pt,
  label={\color{soft}\textbullet}]\small}{\end{itemize}}

% Ligne de titre : intitulé en gras, précision grise
\newcommand{\head}[2]{{\bfseries #1}\if\relax\detokenize{#2}\relax\else{\small\color{soft}\hspace{0.5em}·\hspace{0.5em}#2}\fi\par}
% Ligne de liens d'un projet (démo, code)
\newcommand{\links}[1]{{\footnotesize\color{soft}#1}\par}
\newcommand{\tech}[1]{{\footnotesize\color{soft}#1}\par}
\newcommand{\sep}{\hspace{0.5em}{\color{hair}|}\hspace{0.5em}}

\begin{document}

% =====================================================================
%  EN-TÊTE
% =====================================================================
\begin{minipage}[b]{\textwidth-2.6cm}
  {\fontsize{27}{29}\selectfont\bfseries Ibrahim Kakaev}\par\vspace{5pt}
  {\large Développeur full-stack Java / Spring Boot · Angular\hspace{0.45em}·\hspace{0.45em}Stage de fin d'études 6 mois}\par\vspace{3pt}
  {\small\bfseries À partir de janvier 2027\hspace{0.45em}·\hspace{0.45em}Disponible immédiatement, début anticipé possible}\par\vspace{6pt}
  {\small
   \href{mailto:ibrahim.kakaev@gmail.com}{ibrahim.kakaev@gmail.com}\sep 06 38 05 40 02\sep Albi / Toulouse\sep Nationalité française\sep Permis B\par\vspace{2.5pt}
   \href{https://linkedin.com/in/ibrahim-kakaev-7b11282a0}{LinkedIn}\sep
   \href{https://github.com/IbrahimKakaev}{GitHub}\sep
   \href{https://ibrahimkakaev.github.io/}{Portfolio}\sep
   \href{https://www.root-me.org/MIK_Sparrow?lang=fr}{Root-Me}}
\end{minipage}%
\hfill%
\begin{minipage}[b]{2.35cm}\raggedleft
  \IfFileExists{photo.PNG}{\includegraphics[width=2.35cm,height=2.75cm,keepaspectratio]{photo.PNG}}{}
\end{minipage}

% =====================================================================
\cvsection{Profil}
{\small Master 2 Ingénierie Logicielle. Je construis des applications complètes, de l'API Spring Boot ou Rust à l'interface Angular, livrées testées : intégration continue, analyse de sécurité et démo en ligne. Intérêt marqué pour la sécurité applicative.}\par

% =====================================================================
\cvsection{Compétences techniques}
\begin{grid}\setlength{\itemsep}{1.5pt}
  \item[Langages]      {\small Java 21, TypeScript, JavaScript, Rust, Python, C++, PHP, SQL}
  \item[Back-end]      {\small Spring Boot, Spring Security, JPA, Flyway, Jakarta EE, Node.js, Express.js, API REST, WebSocket, SSE}
  \item[Front-end]     {\small Angular (signals, composants autonomes), RxJS, React}
  \item[Données]       {\small PostgreSQL, pgvector, MongoDB, MySQL, SQLite}
  \item[DevOps]        {\small Linux, Git, Docker, Docker Compose, GitHub Actions, CodeQL, Trivy, AWS}
  \item[Tests, méthodes]{\small JUnit, Testcontainers, Jest, Vitest, PyTest\sep UML, MVC, Scrum, Jira}
  \item[IA et sécurité]{\small RAG, embeddings (ONNX), API LLM ; JWT, OWASP, pentest (Nmap, Burp Suite, Metasploit)}
\end{grid}

% =====================================================================
\cvsection{Projets}
\begin{grid}
  \item[2026]
    \head{Sésame, recherche documentaire d'entreprise par IA}{Java, Spring Boot, pgvector, Angular}
    \begin{cvitems}
      \item Un LLM répond sur des PDF en citant ses sources, sans jamais recevoir un passage interdit à l'utilisateur : filtrage par groupes dans la requête SQL, avant tout appel au modèle.
      \item Tests d'étanchéité de chaque couple utilisateur-document sur chaque route, tests de fuite sur la requête envoyée au LLM, révocation immédiate ; modèle de menace documenté.
      \item JWT avec rotation et détection de rejeu, journal d'audit ; 24 sources sur 24 retrouvées à l'évaluation ; CI avec CodeQL et Trivy.
    \end{cvitems}
    \links{\href{https://ibrahimkakaev.github.io/Sesame/}{Démo : ibrahimkakaev.github.io/Sesame}\sep \href{https://github.com/IbrahimKakaev/Sesame}{Code : github.com/IbrahimKakaev/Sesame}}
  \item[2026]
    \head{NetWatch IDS, détection d'intrusion réseau en temps réel}{Rust, axum, Angular}
    \begin{cvitems}
      \item Moteur de capture (libpcap) et de détection en Rust : décodage Ethernet, VLAN, IPv4/IPv6, TCP, UDP ; détection de scans de ports, SYN flood et débit anormal sur fenêtre glissante, avec mémoire bornée.
      \item Dashboard Angular en WebSocket (carte du trafic en direct), historique SQLite ; tests Rust et Angular en CI.
    \end{cvitems}
    \links{\href{https://ibrahimkakaev.github.io/NetWatch-IDS/}{Démo : ibrahimkakaev.github.io/NetWatch-IDS}\sep \href{https://github.com/IbrahimKakaev/NetWatch-IDS}{Code : github.com/IbrahimKakaev/NetWatch-IDS}}
  \item[2026]
    \head{CERISoNet, réseau social en temps réel}{Angular, Node.js, Socket.IO, MongoDB}
    \begin{cvitems}
      \item Fil d'actualité, likes, commentaires et présence en temps réel ; PostgreSQL pour les comptes, MongoDB pour les publications.
      \item 269 tests (API, sécurité, temps réel), CodeQL et Lighthouse en CI, lancement en une commande avec Docker Compose.
    \end{cvitems}
    \links{\href{https://ibrahimkakaev.github.io/CERISoNet/}{Démo : ibrahimkakaev.github.io/CERISoNet}\sep \href{https://github.com/IbrahimKakaev/CERISoNet}{Code : github.com/IbrahimKakaev/CERISoNet}}
  \item[2025]
    \head{Pentest et cybersécurité offensive}{Kali Linux, Burp Suite, Metasploit}
    \begin{cvitems}
      \item 5 machines VulnHub compromises, 7 modules TryHackMe et 10 challenges Root-Me ; 13 rapports rédigés (reconnaissance, exploitation, élévation de privilèges, remédiation).
    \end{cvitems}
    \links{\href{https://ibrahimkakaev.github.io/reports.html}{Rapports : ibrahimkakaev.github.io/reports.html}\sep \href{https://www.root-me.org/MIK_Sparrow?lang=fr}{Root-Me : MIK\_Sparrow}}
\end{grid}

% =====================================================================
\cvsection{Expérience professionnelle}
\begin{grid}
  \item[2024 · 2 mois]
    \head{Développeur full-stack, stage}{DSI Grand Albigeois, Albi}
    \begin{cvitems}
      \item Interface web de recherche pour les secrétaires d'accueil de la mairie, utilisée en production.
      \item Structuration de données issues de documents scannés et conception de la base de données.
    \end{cvitems}
    \tech{PHP, JavaScript, MySQL}
\end{grid}

% =====================================================================
\cvsection{Formation}
\begin{grid}\setlength{\itemsep}{2pt}
  \item[2024 -- 2027] \head{Master Informatique, Ingénierie Logicielle}{Avignon Université}
  \item[2021 -- 2024] \head{Licence Informatique, option Technologie du Web}{INU Champollion, Albi}
  \item[2019 -- 2021] \head{Baccalauréat, spécialités Informatique et Anglais}{Lycée Lapérouse, Albi}
\end{grid}

% =====================================================================
\cvsection{Autres expériences}
\begin{grid}\setlength{\itemsep}{2pt}
  \item[2020 -- 2024] \head{Animateur jeunesse et référent d'équipe}{CROUS, Grand Albigeois, Handiligue}
  \item[2024 -- 2026] \head{Employé polyvalent et logistique, intérim}{Action, Leclerc Drive, Basic-Fit, Domino's}
\end{grid}

% =====================================================================
\cvsection{Langues et centres d'intérêt}
\begin{grid}\setlength{\itemsep}{1.5pt}
  \item[Langues]   {\small Français langue maternelle\sep Anglais B2\sep Russe courant}
  \item[Intérêts]  {\small Automatisation, musculation, randonnée, musique}
\end{grid}

\end{document}
